\documentclass[11pt,a4paper]{article}
\usepackage[utf8]{inputenc}
\usepackage{amsmath,amssymb}
\usepackage{graphicx}
\usepackage{hyperref}
\usepackage[numbers,sort&compress]{natbib}

\title{Boundary-Aware PINNsformer: Observation-Driven Recovery of Transboundary Pollutant Inflow using Advection-Aligned Pseudo-Sequences}
\author{BBM Research}
\date{\today}

\begin{document}
\maketitle

\begin{abstract}
City-scale air quality models must specify what enters the domain across lateral boundaries.
We invert the problem: given sparse interior CPCB station observations and a known ERA5 wind field over the Delhi-NCR airshed, we jointly recover the interior concentration field, the time-varying direction-resolved boundary inflow, and the interior source field subject to the advection--diffusion--deposition PDE as a soft constraint, using a PINNsformer whose pseudo-sequences are sampled along backward wind characteristics.
\end{abstract}

\section{Introduction}
% PRD §11: policy problem + inventory-free framing. RQ1–RQ4; transboundary share as headline number.
\begin{itemize}
  \item Delhi-NCR winter/post-monsoon burden; why lateral boundary inflow is the missing input.
  \item Inversion framing: boundary inflow as the unknown, not an input (contrast nested CTMs).
  \item Contributions: inverse-boundary formulation; first PINNsformer on real observations; advection-aligned pseudo-sequences; identifiability boundary; inventory-free share validated on withheld data.
\end{itemize}

\section{Related work}
% PRD §11 + notes/gap_matrix.md: three-literature gap matrix. Every claim verified against primary sources.
\begin{itemize}
  \item PINNsformer line and successors; pseudo-sequence idea and its limits on benchmark PDEs.
  \item Physics-informed air quality (AirPhyNet, SPIN, CoNOAir, OmniAir): boundary treatment gap.
  \item PINN inverse-boundary work (acoustics, haemodynamics): no atmospheric transboundary analogue.
  \item Attribution practice: DSS, WRF-Chem inventory-based shares as external reference, not baseline.
\end{itemize}

\section{Method}
% PRD §2: contracts — domain/arc-length s, nets (C, C_b, S, K, lam), 4 pseudo-sequence variants, wind-switched BCs, GradNorm-balanced objective, curriculum.
\begin{itemize}
  \item Governing equation $\partial C/\partial t + \nabla\cdot(\mathbf{u}C) = \nabla\cdot(K\nabla C) + S - \lambda C$; 2-D vertically integrated; metric CRS; normalizer owns scaling.
  \item Networks: $C$ (PINNsformer over $L$ tokens), $C_b(s,t)$ (softplus, temporal smoothness), $S(x,y,t)$ (non-negative, L1/TV sparsity), $K$, $\lambda$ (log-params; PM1 $<$ PM2.5 $<$ PM10 ordering).
  \item Pseudo-sequences: uniform\_forward / uniform\_backward / advection\_backward (ours) / advection\_jitter (ours); RK2/RK4 on interpolated ERA5 wind; boundary-hit flags.
  \item Boundary switching: Dirichlet $C - C_b$ on inflow ($\mathbf{u}\cdot\hat{\mathbf{n}} < 0$), zero-gradient on outflow, hysteresis band.
  \item Objective $L = w_{\mathrm{data}}L_{\mathrm{data}} + w_{\mathrm{pde}}L_{\mathrm{pde}} + w_{\mathrm{bc}}L_{\mathrm{bc}} + w_{\mathrm{reg}}L_{\mathrm{reg}}$; masked data loss, gradient-norm balancing, Adam $\to$ L-BFGS, temporal curriculum.
  \item Identifiability argument: wind directional entropy as tomographic angular diversity.
\end{itemize}

\section{Data}
% PRD E0: CPCB CCR + OpenCity, ERA5, withheld FIRMS; QC/mask (never impute); domain + interior/perimeter split; Oct 2019–Feb 2024.
\begin{itemize}
  \item CPCB CCR + OpenCity mirror: pollutants, met; harmonization; station inclusion threshold.
  \item ERA5 (CDS): u10/v10, BLH, T, RH; caching; station-vs-ERA5 wind agreement.
  \item QC: range, flatline, spike, duplicates, IST, units; validity masks; missingness characterization.
  \item Domain: Delhi-NCR box, metric CRS, arc-length $s$, interior vs NCR perimeter ring; split hashes.
  \item FIRMS VIIRS/MODIS fire + FRP: held out entirely (leakage guard); E4 use only.
\end{itemize}

\section{Experiments}
% PRD §6: E0–E5 with acceptance criteria; §7 baselines + metrics.
\begin{itemize}
  \item E0 data pipeline + characterization (one-command rebuild; wind agreement; entropy per window).
  \item E1 synthetic identifiability: FD ground truth; sweep stations $\{7,15,25,40\}$ $\times$ config $\times$ entropy $\times$ noise; RQ2.
  \item E2 real recovery: interior fit, withheld perimeter validation; per-season/pollutant; RQ1.
  \item E3 ablations: 4 pseudo-sequence variants at matched params; $L$ sweep; Wavelet; physics term; RQ3.
  \item E4 fire validation: NW-sector inflow vs withheld FRP (Spearman + lags); DSS/WRF-Chem comparison; RQ4.
  \item E5 transfer: Kolkata/Chennai domain override.
  \item Baselines: zero/climatological inflow, MLP-PINN, uniform-PSF PINNsformer, back-trajectory regression, published inventory shares (external).
\end{itemize}

\section{Results}
% Figures/tables generated from results/ via scripts 90/91; claims ledger maps each claim to an artifact.
\begin{itemize}
  \item E1 identifiability surface (error vs stations $\times$ entropy); tomography hypothesis test.
  \item E2 perimeter RMSE/MAE/R$^2$ vs baselines; $C_b(s,t)$ Hovm\"oller by season; transboundary share.
  \item E3 accuracy vs $L$ / wall-clock / memory per variant; Wavelet + physics-term verdicts.
  \item E4 NW inflow vs FRP correlation; external share comparison and divergence narrative.
  \item E5 second-airshed outcome; physical consistency (mass balance; $K$, $\lambda$ ranges; deposition ordering).
\end{itemize}

\section{Discussion}
% PRD §11: what the transboundary number means; where/why it diverges from DSS/WRF-Chem.
\begin{itemize}
  \item Reading of the inventory-free share per season/pollutant; sector (NW/N/E/S/W) decomposition.
  \item Divergence vs DSS/WRF-Chem: accounting differences, inventory vintage, vertical-mixing and wind-error contributions.
  \item Attention attribution: boundary-hitting tokens as visual justification of advection alignment.
\end{itemize}

\section{Limitations}
% PRD §11 + risk register: 2-D vertical integration, winter inversions, wind error, network sparsity.
\begin{itemize}
  \item 2-D vertically integrated field; BLH-modulated $K$/loss only partly captures winter inversions.
  \item ERA5 wind error propagation; claims restricted to good-agreement periods; perturbed-wind sensitivity.
  \item Sparse/perimeter station geometry bounds E2 validation strength; E1 error floor applies.
  \item Single-species first-order loss; no chemistry/aerosol microphysics; $S$ is a nuisance field, not an inventory.
  \item No forecasting claim.
\end{itemize}

\section{Conclusion}
% Target venues: Environmental Modelling \& Software (primary), Atmospheric Environment, GMD, SIGSPATIAL/AAAI AISI, Sci.\ Reports/IEEE Access (fallback).
\begin{itemize}
  \item Summary of RQ1–RQ4 answers; fallback: E1 identifiability boundary alone is publishable.
  \item Release: code + processed dataset + configs reproducing every figure.
\end{itemize}


\bibliographystyle{plainnat}
\bibliography{refs}

\end{document}
